

\subsection{Create a horizontal Grid: \texttt{gridCreate}}
\index{gridCreate}
\label{gridCreate}

The function {\texttt{gridCreate}} creates a horizontal Grid.

\subsubsection*{Usage}

\begin{verbatim}
    INTEGER FUNCTION gridCreate(INTEGER gridtype, INTEGER size)
\end{verbatim}

\hspace*{4mm}\begin{minipage}[]{15cm}
\begin{deflist}{\texttt{gridtype}\ }
\item[\texttt{gridtype}]
The type of the grid, one of the set of predefined {\CDI} grid types.
                     The valid {\CDI} grid types are {\texttt{GRID\_GENERIC}}, {\texttt{GRID\_GAUSSIAN}},
                     {\texttt{GRID\_LONLAT}}, {\texttt{GRID\_PROJECTION}}, {\texttt{GRID\_SPECTRAL}},
                     {\texttt{GRID\_GME}}, {\texttt{GRID\_CURVILINEAR}} and {\texttt{GRID\_UNSTRUCTURED}}.
\item[\texttt{size}]
Number of gridpoints.

\end{deflist}
\end{minipage}

\subsubsection*{Result}

{\texttt{gridCreate}} returns an identifier to the Grid.


\subsubsection*{Example}

Here is an example using {\texttt{gridCreate}} to create a regular lon/lat Grid:

\begin{lstlisting}[language=Fortran, backgroundcolor=\color{pyellow}, basicstyle=\small, columns=flexible]

    INCLUDE 'cdi.h'
       ...
    #define  nlon  12
    #define  nlat   6
       ...
    REAL*8 lons(nlon) = (/0, 30, 60, 90, 120, 150, 180, 210, 240, 270, 300, 330/)
    REAL*8 lats(nlat) = (/-75, -45, -15, 15, 45, 75/)
    INTEGER gridID
       ...
    gridID = gridCreate(GRID_LONLAT, nlon*nlat)
    CALL gridDefXsize(gridID, nlon)
    CALL gridDefYsize(gridID, nlat)
    CALL gridDefXvals(gridID, lons)
    CALL gridDefYvals(gridID, lats)
       ...
\end{lstlisting}


\subsection{Destroy a horizontal Grid: \texttt{gridDestroy}}
\index{gridDestroy}
\label{gridDestroy}
\subsubsection*{Usage}

\begin{verbatim}
    SUBROUTINE gridDestroy(INTEGER gridID)
\end{verbatim}

\hspace*{4mm}\begin{minipage}[]{15cm}
\begin{deflist}{\texttt{gridID}\ }
\item[\texttt{gridID}]
Grid ID, from a previous call to {\htmlref{\texttt{gridCreate}}{gridCreate}}.

\end{deflist}
\end{minipage}


\subsection{Duplicate a horizontal Grid: \texttt{gridDuplicate}}
\index{gridDuplicate}
\label{gridDuplicate}

The function {\texttt{gridDuplicate}} duplicates a horizontal Grid.

\subsubsection*{Usage}

\begin{verbatim}
    INTEGER FUNCTION gridDuplicate(INTEGER gridID)
\end{verbatim}

\hspace*{4mm}\begin{minipage}[]{15cm}
\begin{deflist}{\texttt{gridID}\ }
\item[\texttt{gridID}]
Grid ID, from a previous call to {\htmlref{\texttt{gridCreate}}{gridCreate}} or {\htmlref{\texttt{vlistInqVarGrid}}{vlistInqVarGrid}}.

\end{deflist}
\end{minipage}

\subsubsection*{Result}

{\texttt{gridDuplicate}} returns an identifier to the duplicated Grid.



\subsection{Get the type of a Grid: \texttt{gridInqType}}
\index{gridInqType}
\label{gridInqType}

The function {\texttt{gridInqType}} returns the type of a Grid.

\subsubsection*{Usage}

\begin{verbatim}
    INTEGER FUNCTION gridInqType(INTEGER gridID)
\end{verbatim}

\hspace*{4mm}\begin{minipage}[]{15cm}
\begin{deflist}{\texttt{gridID}\ }
\item[\texttt{gridID}]
Grid ID, from a previous call to {\htmlref{\texttt{gridCreate}}{gridCreate}} or {\htmlref{\texttt{vlistInqVarGrid}}{vlistInqVarGrid}}.

\end{deflist}
\end{minipage}

\subsubsection*{Result}

{\texttt{gridInqType}} returns the type of the grid,
one of the set of predefined {\CDI} grid types.
The valid {\CDI} grid types are {\texttt{GRID\_GENERIC}}, {\texttt{GRID\_GAUSSIAN}},
{\texttt{GRID\_LONLAT}}, {\texttt{GRID\_PROJECTION}}, {\texttt{GRID\_SPECTRAL}}, {\texttt{GRID\_GME}},
{\texttt{GRID\_CURVILINEAR}} and {\texttt{GRID\_UNSTRUCTURED}}.



\subsection{Get the size of a Grid: \texttt{gridInqSize}}
\index{gridInqSize}
\label{gridInqSize}

The function {\texttt{gridInqSize}} returns the size of a Grid.

\subsubsection*{Usage}

\begin{verbatim}
    INTEGER FUNCTION gridInqSize(INTEGER gridID)
\end{verbatim}

\hspace*{4mm}\begin{minipage}[]{15cm}
\begin{deflist}{\texttt{gridID}\ }
\item[\texttt{gridID}]
Grid ID, from a previous call to {\htmlref{\texttt{gridCreate}}{gridCreate}} or {\htmlref{\texttt{vlistInqVarGrid}}{vlistInqVarGrid}}.

\end{deflist}
\end{minipage}

\subsubsection*{Result}

{\texttt{gridInqSize}} returns the number of grid points of a Grid.



\subsection{Define the number of values of a X-axis: \texttt{gridDefXsize}}
\index{gridDefXsize}
\label{gridDefXsize}

The function {\texttt{gridDefXsize}} defines the number of values of a X-axis.

\subsubsection*{Usage}

\begin{verbatim}
    SUBROUTINE gridDefXsize(INTEGER gridID, INTEGER xsize)
\end{verbatim}

\hspace*{4mm}\begin{minipage}[]{15cm}
\begin{deflist}{\texttt{gridID}\ }
\item[\texttt{gridID}]
Grid ID, from a previous call to {\htmlref{\texttt{gridCreate}}{gridCreate}}.
\item[\texttt{xsize}]
Number of values of a X-axis.

\end{deflist}
\end{minipage}


\subsection{Get the number of values of a X-axis: \texttt{gridInqXsize}}
\index{gridInqXsize}
\label{gridInqXsize}

The function {\texttt{gridInqXsize}} returns the number of values of a X-axis.

\subsubsection*{Usage}

\begin{verbatim}
    INTEGER FUNCTION gridInqXsize(INTEGER gridID)
\end{verbatim}

\hspace*{4mm}\begin{minipage}[]{15cm}
\begin{deflist}{\texttt{gridID}\ }
\item[\texttt{gridID}]
Grid ID, from a previous call to {\htmlref{\texttt{gridCreate}}{gridCreate}} or {\htmlref{\texttt{vlistInqVarGrid}}{vlistInqVarGrid}}.

\end{deflist}
\end{minipage}

\subsubsection*{Result}

{\texttt{gridInqXsize}} returns the number of values of a X-axis.



\subsection{Define the number of values of a Y-axis: \texttt{gridDefYsize}}
\index{gridDefYsize}
\label{gridDefYsize}

The function {\texttt{gridDefYsize}} defines the number of values of a Y-axis.

\subsubsection*{Usage}

\begin{verbatim}
    SUBROUTINE gridDefYsize(INTEGER gridID, INTEGER ysize)
\end{verbatim}

\hspace*{4mm}\begin{minipage}[]{15cm}
\begin{deflist}{\texttt{gridID}\ }
\item[\texttt{gridID}]
Grid ID, from a previous call to {\htmlref{\texttt{gridCreate}}{gridCreate}}.
\item[\texttt{ysize}]
Number of values of a Y-axis.

\end{deflist}
\end{minipage}


\subsection{Get the number of values of a Y-axis: \texttt{gridInqYsize}}
\index{gridInqYsize}
\label{gridInqYsize}

The function {\texttt{gridInqYsize}} returns the number of values of a Y-axis.

\subsubsection*{Usage}

\begin{verbatim}
    INTEGER FUNCTION gridInqYsize(INTEGER gridID)
\end{verbatim}

\hspace*{4mm}\begin{minipage}[]{15cm}
\begin{deflist}{\texttt{gridID}\ }
\item[\texttt{gridID}]
Grid ID, from a previous call to {\htmlref{\texttt{gridCreate}}{gridCreate}} or {\htmlref{\texttt{vlistInqVarGrid}}{vlistInqVarGrid}}.

\end{deflist}
\end{minipage}

\subsubsection*{Result}

{\texttt{gridInqYsize}} returns the number of values of a Y-axis.



\subsection{Define the number of parallels between a pole and the equator: \texttt{gridDefNP}}
\index{gridDefNP}
\label{gridDefNP}

The function {\texttt{gridDefNP}} defines the number of parallels between a pole and the equator
of a Gaussian grid.

\subsubsection*{Usage}

\begin{verbatim}
    SUBROUTINE gridDefNP(INTEGER gridID, INTEGER np)
\end{verbatim}

\hspace*{4mm}\begin{minipage}[]{15cm}
\begin{deflist}{\texttt{gridID}\ }
\item[\texttt{gridID}]
Grid ID, from a previous call to {\htmlref{\texttt{gridCreate}}{gridCreate}}.
\item[\texttt{np}]
Number of parallels between a pole and the equator.

\end{deflist}
\end{minipage}


\subsection{Get the number of parallels between a pole and the equator: \texttt{gridInqNP}}
\index{gridInqNP}
\label{gridInqNP}

The function {\texttt{gridInqNP}} returns the number of parallels between a pole and the equator
of a Gaussian grid.

\subsubsection*{Usage}

\begin{verbatim}
    INTEGER FUNCTION gridInqNP(INTEGER gridID)
\end{verbatim}

\hspace*{4mm}\begin{minipage}[]{15cm}
\begin{deflist}{\texttt{gridID}\ }
\item[\texttt{gridID}]
Grid ID, from a previous call to {\htmlref{\texttt{gridCreate}}{gridCreate}} or {\htmlref{\texttt{vlistInqVarGrid}}{vlistInqVarGrid}}.

\end{deflist}
\end{minipage}

\subsubsection*{Result}

{\texttt{gridInqNP}} returns the number of parallels between a pole and the equator.



\subsection{Define the values of a X-axis: \texttt{gridDefXvals}}
\index{gridDefXvals}
\label{gridDefXvals}

The function {\texttt{gridDefXvals}} defines all values of the X-axis.

\subsubsection*{Usage}

\begin{verbatim}
    SUBROUTINE gridDefXvals(INTEGER gridID, REAL*8 xvals)
\end{verbatim}

\hspace*{4mm}\begin{minipage}[]{15cm}
\begin{deflist}{\texttt{gridID}\ }
\item[\texttt{gridID}]
Grid ID, from a previous call to {\htmlref{\texttt{gridCreate}}{gridCreate}}.
\item[\texttt{xvals}]
X-values of the grid.

\end{deflist}
\end{minipage}


\subsection{Get all values of a X-axis: \texttt{gridInqXvals}}
\index{gridInqXvals}
\label{gridInqXvals}

The function {\texttt{gridInqXvals}} returns all values of the X-axis.

\subsubsection*{Usage}

\begin{verbatim}
    INTEGER FUNCTION gridInqXvals(INTEGER gridID, REAL*8 xvals)
\end{verbatim}

\hspace*{4mm}\begin{minipage}[]{15cm}
\begin{deflist}{\texttt{gridID}\ }
\item[\texttt{gridID}]
Grid ID, from a previous call to {\htmlref{\texttt{gridCreate}}{gridCreate}} or {\htmlref{\texttt{vlistInqVarGrid}}{vlistInqVarGrid}}.
\item[\texttt{xvals}]
Pointer to the location into which the X-values are read.
                    The caller must allocate space for the returned values.

\end{deflist}
\end{minipage}

\subsubsection*{Result}

Upon successful completion {\texttt{gridInqXvals}} returns the number of values and
the values are stored in {\texttt{xvals}}.
Otherwise, 0 is returned and {\texttt{xvals}} is empty.



\subsection{Define the values of a Y-axis: \texttt{gridDefYvals}}
\index{gridDefYvals}
\label{gridDefYvals}

The function {\texttt{gridDefYvals}} defines all values of the Y-axis.

\subsubsection*{Usage}

\begin{verbatim}
    SUBROUTINE gridDefYvals(INTEGER gridID, REAL*8 yvals)
\end{verbatim}

\hspace*{4mm}\begin{minipage}[]{15cm}
\begin{deflist}{\texttt{gridID}\ }
\item[\texttt{gridID}]
Grid ID, from a previous call to {\htmlref{\texttt{gridCreate}}{gridCreate}}.
\item[\texttt{yvals}]
Y-values of the grid.

\end{deflist}
\end{minipage}


\subsection{Get all values of a Y-axis: \texttt{gridInqYvals}}
\index{gridInqYvals}
\label{gridInqYvals}

The function {\texttt{gridInqYvals}} returns all values of the Y-axis.

\subsubsection*{Usage}

\begin{verbatim}
    INTEGER FUNCTION gridInqYvals(INTEGER gridID, REAL*8 yvals)
\end{verbatim}

\hspace*{4mm}\begin{minipage}[]{15cm}
\begin{deflist}{\texttt{gridID}\ }
\item[\texttt{gridID}]
Grid ID, from a previous call to {\htmlref{\texttt{gridCreate}}{gridCreate}} or {\htmlref{\texttt{vlistInqVarGrid}}{vlistInqVarGrid}}.
\item[\texttt{yvals}]
Pointer to the location into which the Y-values are read.
                    The caller must allocate space for the returned values.

\end{deflist}
\end{minipage}

\subsubsection*{Result}

Upon successful completion {\texttt{gridInqYvals}} returns the number of values and
the values are stored in {\texttt{yvals}}.
Otherwise, 0 is returned and {\texttt{yvals}} is empty.



\subsection{Define the bounds of a X-axis: \texttt{gridDefXbounds}}
\index{gridDefXbounds}
\label{gridDefXbounds}

The function {\texttt{gridDefXbounds}} defines all bounds of the X-axis.

\subsubsection*{Usage}

\begin{verbatim}
    SUBROUTINE gridDefXbounds(INTEGER gridID, REAL*8 xbounds)
\end{verbatim}

\hspace*{4mm}\begin{minipage}[]{15cm}
\begin{deflist}{\texttt{xbounds}\ }
\item[\texttt{gridID}]
Grid ID, from a previous call to {\htmlref{\texttt{gridCreate}}{gridCreate}}.
\item[\texttt{xbounds}]
X-bounds of the grid.

\end{deflist}
\end{minipage}


\subsection{Get the bounds of a X-axis: \texttt{gridInqXbounds}}
\index{gridInqXbounds}
\label{gridInqXbounds}

The function {\texttt{gridInqXbounds}} returns the bounds of the X-axis.

\subsubsection*{Usage}

\begin{verbatim}
    INTEGER FUNCTION gridInqXbounds(INTEGER gridID, REAL*8 xbounds)
\end{verbatim}

\hspace*{4mm}\begin{minipage}[]{15cm}
\begin{deflist}{\texttt{xbounds}\ }
\item[\texttt{gridID}]
Grid ID, from a previous call to {\htmlref{\texttt{gridCreate}}{gridCreate}} or {\htmlref{\texttt{vlistInqVarGrid}}{vlistInqVarGrid}}.
\item[\texttt{xbounds}]
Pointer to the location into which the X-bounds are read.
                    The caller must allocate space for the returned values.

\end{deflist}
\end{minipage}

\subsubsection*{Result}

Upon successful completion {\texttt{gridInqXbounds}} returns the number of bounds and
the bounds are stored in {\texttt{xbounds}}.
Otherwise, 0 is returned and {\texttt{xbounds}} is empty.



\subsection{Define the bounds of a Y-axis: \texttt{gridDefYbounds}}
\index{gridDefYbounds}
\label{gridDefYbounds}

The function {\texttt{gridDefYbounds}} defines all bounds of the Y-axis.

\subsubsection*{Usage}

\begin{verbatim}
    SUBROUTINE gridDefYbounds(INTEGER gridID, REAL*8 ybounds)
\end{verbatim}

\hspace*{4mm}\begin{minipage}[]{15cm}
\begin{deflist}{\texttt{ybounds}\ }
\item[\texttt{gridID}]
Grid ID, from a previous call to {\htmlref{\texttt{gridCreate}}{gridCreate}}.
\item[\texttt{ybounds}]
Y-bounds of the grid.

\end{deflist}
\end{minipage}


\subsection{Get the bounds of a Y-axis: \texttt{gridInqYbounds}}
\index{gridInqYbounds}
\label{gridInqYbounds}

The function {\texttt{gridInqYbounds}} returns the bounds of the Y-axis.

\subsubsection*{Usage}

\begin{verbatim}
    INTEGER FUNCTION gridInqYbounds(INTEGER gridID, REAL*8 ybounds)
\end{verbatim}

\hspace*{4mm}\begin{minipage}[]{15cm}
\begin{deflist}{\texttt{ybounds}\ }
\item[\texttt{gridID}]
Grid ID, from a previous call to {\htmlref{\texttt{gridCreate}}{gridCreate}} or {\htmlref{\texttt{vlistInqVarGrid}}{vlistInqVarGrid}}.
\item[\texttt{ybounds}]
Pointer to the location into which the Y-bounds are read.
                    The caller must allocate space for the returned values.

\end{deflist}
\end{minipage}

\subsubsection*{Result}

Upon successful completion {\texttt{gridInqYbounds}} returns the number of bounds and
the bounds are stored in {\texttt{ybounds}}.
Otherwise, 0 is returned and {\texttt{ybounds}} is empty.



\subsection{Define the name of a X-axis: \texttt{gridDefXname}}
\index{gridDefXname}
\label{gridDefXname}

The function {\texttt{gridDefXname}} defines the name of a X-axis.

\subsubsection*{Usage}

\begin{verbatim}
    SUBROUTINE gridDefXname(INTEGER gridID, CHARACTER*(*) name)
\end{verbatim}

\hspace*{4mm}\begin{minipage}[]{15cm}
\begin{deflist}{\texttt{gridID}\ }
\item[\texttt{gridID}]
Grid ID, from a previous call to {\htmlref{\texttt{gridCreate}}{gridCreate}}.
\item[\texttt{name}]
Name of the X-axis.

\end{deflist}
\end{minipage}


\subsection{Get the name of a X-axis: \texttt{gridInqXname}}
\index{gridInqXname}
\label{gridInqXname}

The function {\texttt{gridInqXname}} returns the name of a X-axis.

\subsubsection*{Usage}

\begin{verbatim}
    SUBROUTINE gridInqXname(INTEGER gridID, CHARACTER*(*) name)
\end{verbatim}

\hspace*{4mm}\begin{minipage}[]{15cm}
\begin{deflist}{\texttt{gridID}\ }
\item[\texttt{gridID}]
Grid ID, from a previous call to {\htmlref{\texttt{gridCreate}}{gridCreate}} or {\htmlref{\texttt{vlistInqVarGrid}}{vlistInqVarGrid}}.
\item[\texttt{name}]
Name of the X-axis. The caller must allocate space for the
                    returned string. The maximum possible length, in characters, of
                    the string is given by the predefined constant {\texttt{CDI\_MAX\_NAME}}.

\end{deflist}
\end{minipage}

\subsubsection*{Result}

{\texttt{gridInqXname}} returns the name of the X-axis to the parameter name.



\subsection{Define the longname of a X-axis: \texttt{gridDefXlongname}}
\index{gridDefXlongname}
\label{gridDefXlongname}

The function {\texttt{gridDefXlongname}} defines the longname of a X-axis.

\subsubsection*{Usage}

\begin{verbatim}
    SUBROUTINE gridDefXlongname(INTEGER gridID, CHARACTER*(*) longname)
\end{verbatim}

\hspace*{4mm}\begin{minipage}[]{15cm}
\begin{deflist}{\texttt{longname}\ }
\item[\texttt{gridID}]
Grid ID, from a previous call to {\htmlref{\texttt{gridCreate}}{gridCreate}}.
\item[\texttt{longname}]
Longname of the X-axis.

\end{deflist}
\end{minipage}


\subsection{Get the longname of a X-axis: \texttt{gridInqXlongname}}
\index{gridInqXlongname}
\label{gridInqXlongname}

The function {\texttt{gridInqXlongname}} returns the longname of a X-axis.

\subsubsection*{Usage}

\begin{verbatim}
    SUBROUTINE gridInqXlongname(INTEGER gridID, CHARACTER*(*) longname)
\end{verbatim}

\hspace*{4mm}\begin{minipage}[]{15cm}
\begin{deflist}{\texttt{longname}\ }
\item[\texttt{gridID}]
Grid ID, from a previous call to {\htmlref{\texttt{gridCreate}}{gridCreate}} or {\htmlref{\texttt{vlistInqVarGrid}}{vlistInqVarGrid}}.
\item[\texttt{longname}]
Longname of the X-axis. The caller must allocate space for the
                    returned string. The maximum possible length, in characters, of
                    the string is given by the predefined constant {\texttt{CDI\_MAX\_NAME}}.

\end{deflist}
\end{minipage}

\subsubsection*{Result}

{\texttt{gridInqXlongname}} returns the longname of the X-axis to the parameter longname.



\subsection{Define the units of a X-axis: \texttt{gridDefXunits}}
\index{gridDefXunits}
\label{gridDefXunits}

The function {\texttt{gridDefXunits}} defines the units of a X-axis.

\subsubsection*{Usage}

\begin{verbatim}
    SUBROUTINE gridDefXunits(INTEGER gridID, CHARACTER*(*) units)
\end{verbatim}

\hspace*{4mm}\begin{minipage}[]{15cm}
\begin{deflist}{\texttt{gridID}\ }
\item[\texttt{gridID}]
Grid ID, from a previous call to {\htmlref{\texttt{gridCreate}}{gridCreate}}.
\item[\texttt{units}]
Units of the X-axis.

\end{deflist}
\end{minipage}


\subsection{Get the units of a X-axis: \texttt{gridInqXunits}}
\index{gridInqXunits}
\label{gridInqXunits}

The function {\texttt{gridInqXunits}} returns the units of a X-axis.

\subsubsection*{Usage}

\begin{verbatim}
    SUBROUTINE gridInqXunits(INTEGER gridID, CHARACTER*(*) units)
\end{verbatim}

\hspace*{4mm}\begin{minipage}[]{15cm}
\begin{deflist}{\texttt{gridID}\ }
\item[\texttt{gridID}]
Grid ID, from a previous call to {\htmlref{\texttt{gridCreate}}{gridCreate}} or {\htmlref{\texttt{vlistInqVarGrid}}{vlistInqVarGrid}}.
\item[\texttt{units}]
Units of the X-axis. The caller must allocate space for the
                    returned string. The maximum possible length, in characters, of
                    the string is given by the predefined constant {\texttt{CDI\_MAX\_NAME}}.

\end{deflist}
\end{minipage}

\subsubsection*{Result}

{\texttt{gridInqXunits}} returns the units of the X-axis to the parameter units.



\subsection{Define the name of a Y-axis: \texttt{gridDefYname}}
\index{gridDefYname}
\label{gridDefYname}

The function {\texttt{gridDefYname}} defines the name of a Y-axis.

\subsubsection*{Usage}

\begin{verbatim}
    SUBROUTINE gridDefYname(INTEGER gridID, CHARACTER*(*) name)
\end{verbatim}

\hspace*{4mm}\begin{minipage}[]{15cm}
\begin{deflist}{\texttt{gridID}\ }
\item[\texttt{gridID}]
Grid ID, from a previous call to {\htmlref{\texttt{gridCreate}}{gridCreate}}.
\item[\texttt{name}]
Name of the Y-axis.

\end{deflist}
\end{minipage}


\subsection{Get the name of a Y-axis: \texttt{gridInqYname}}
\index{gridInqYname}
\label{gridInqYname}

The function {\texttt{gridInqYname}} returns the name of a Y-axis.

\subsubsection*{Usage}

\begin{verbatim}
    SUBROUTINE gridInqYname(INTEGER gridID, CHARACTER*(*) name)
\end{verbatim}

\hspace*{4mm}\begin{minipage}[]{15cm}
\begin{deflist}{\texttt{gridID}\ }
\item[\texttt{gridID}]
Grid ID, from a previous call to {\htmlref{\texttt{gridCreate}}{gridCreate}} or {\htmlref{\texttt{vlistInqVarGrid}}{vlistInqVarGrid}}.
\item[\texttt{name}]
Name of the Y-axis. The caller must allocate space for the
                    returned string. The maximum possible length, in characters, of
                    the string is given by the predefined constant {\texttt{CDI\_MAX\_NAME}}.

\end{deflist}
\end{minipage}

\subsubsection*{Result}

{\texttt{gridInqYname}} returns the name of the Y-axis to the parameter name.



\subsection{Define the longname of a Y-axis: \texttt{gridDefYlongname}}
\index{gridDefYlongname}
\label{gridDefYlongname}

The function {\texttt{gridDefYlongname}} defines the longname of a Y-axis.

\subsubsection*{Usage}

\begin{verbatim}
    SUBROUTINE gridDefYlongname(INTEGER gridID, CHARACTER*(*) longname)
\end{verbatim}

\hspace*{4mm}\begin{minipage}[]{15cm}
\begin{deflist}{\texttt{longname}\ }
\item[\texttt{gridID}]
Grid ID, from a previous call to {\htmlref{\texttt{gridCreate}}{gridCreate}}.
\item[\texttt{longname}]
Longname of the Y-axis.

\end{deflist}
\end{minipage}


\subsection{Get the longname of a Y-axis: \texttt{gridInqYlongname}}
\index{gridInqYlongname}
\label{gridInqYlongname}

The function {\texttt{gridInqYlongname}} returns the longname of a Y-axis.

\subsubsection*{Usage}

\begin{verbatim}
    SUBROUTINE gridInqYlongname(INTEGER gridID, CHARACTER*(*) longname)
\end{verbatim}

\hspace*{4mm}\begin{minipage}[]{15cm}
\begin{deflist}{\texttt{longname}\ }
\item[\texttt{gridID}]
Grid ID, from a previous call to {\htmlref{\texttt{gridCreate}}{gridCreate}} or {\htmlref{\texttt{vlistInqVarGrid}}{vlistInqVarGrid}}.
\item[\texttt{longname}]
Longname of the Y-axis. The caller must allocate space for the
                    returned string. The maximum possible length, in characters, of
                    the string is given by the predefined constant {\texttt{CDI\_MAX\_NAME}}.

\end{deflist}
\end{minipage}

\subsubsection*{Result}

{\texttt{gridInqYlongname}} returns the longname of the Y-axis to the parameter longname.



\subsection{Define the units of a Y-axis: \texttt{gridDefYunits}}
\index{gridDefYunits}
\label{gridDefYunits}

The function {\texttt{gridDefYunits}} defines the units of a Y-axis.

\subsubsection*{Usage}

\begin{verbatim}
    SUBROUTINE gridDefYunits(INTEGER gridID, CHARACTER*(*) units)
\end{verbatim}

\hspace*{4mm}\begin{minipage}[]{15cm}
\begin{deflist}{\texttt{gridID}\ }
\item[\texttt{gridID}]
Grid ID, from a previous call to {\htmlref{\texttt{gridCreate}}{gridCreate}}.
\item[\texttt{units}]
Units of the Y-axis.

\end{deflist}
\end{minipage}


\subsection{Get the units of a Y-axis: \texttt{gridInqYunits}}
\index{gridInqYunits}
\label{gridInqYunits}

The function {\texttt{gridInqYunits}} returns the units of a Y-axis.

\subsubsection*{Usage}

\begin{verbatim}
    SUBROUTINE gridInqYunits(INTEGER gridID, CHARACTER*(*) units)
\end{verbatim}

\hspace*{4mm}\begin{minipage}[]{15cm}
\begin{deflist}{\texttt{gridID}\ }
\item[\texttt{gridID}]
Grid ID, from a previous call to {\htmlref{\texttt{gridCreate}}{gridCreate}} or {\htmlref{\texttt{vlistInqVarGrid}}{vlistInqVarGrid}}.
\item[\texttt{units}]
Units of the Y-axis. The caller must allocate space for the
                    returned string. The maximum possible length, in characters, of
                    the string is given by the predefined constant {\texttt{CDI\_MAX\_NAME}}.

\end{deflist}
\end{minipage}

\subsubsection*{Result}

{\texttt{gridInqYunits}} returns the units of the Y-axis to the parameter units.



\subsection{Define the reference number for an unstructured grid: \texttt{gridDefNumber}}
\index{gridDefNumber}
\label{gridDefNumber}

The function {\texttt{gridDefNumber}} defines the reference number for an unstructured grid.

\subsubsection*{Usage}

\begin{verbatim}
    SUBROUTINE gridDefNumber(INTEGER gridID, INTEGER number)
\end{verbatim}

\hspace*{4mm}\begin{minipage}[]{15cm}
\begin{deflist}{\texttt{gridID}\ }
\item[\texttt{gridID}]
Grid ID, from a previous call to {\htmlref{\texttt{gridCreate}}{gridCreate}}.
\item[\texttt{number}]
Reference number for an unstructured grid.

\end{deflist}
\end{minipage}


\subsection{Get the reference number to an unstructured grid: \texttt{gridInqNumber}}
\index{gridInqNumber}
\label{gridInqNumber}

The function {\texttt{gridInqNumber}} returns the reference number to an unstructured grid.

\subsubsection*{Usage}

\begin{verbatim}
    INTEGER FUNCTION gridInqNumber(INTEGER gridID)
\end{verbatim}

\hspace*{4mm}\begin{minipage}[]{15cm}
\begin{deflist}{\texttt{gridID}\ }
\item[\texttt{gridID}]
Grid ID, from a previous call to {\htmlref{\texttt{gridCreate}}{gridCreate}} or {\htmlref{\texttt{vlistInqVarGrid}}{vlistInqVarGrid}}.

\end{deflist}
\end{minipage}

\subsubsection*{Result}

{\texttt{gridInqNumber}} returns the reference number to an unstructured grid.


\subsection{Define the position of grid in the reference file: \texttt{gridDefPosition}}
\index{gridDefPosition}
\label{gridDefPosition}

The function {\texttt{gridDefPosition}} defines the position of grid in the reference file.

\subsubsection*{Usage}

\begin{verbatim}
    SUBROUTINE gridDefPosition(INTEGER gridID, INTEGER position)
\end{verbatim}

\hspace*{4mm}\begin{minipage}[]{15cm}
\begin{deflist}{\texttt{position}\ }
\item[\texttt{gridID}]
Grid ID, from a previous call to {\htmlref{\texttt{gridCreate}}{gridCreate}}.
\item[\texttt{position}]
Position of grid in the reference file.

\end{deflist}
\end{minipage}


\subsection{Get the position of grid in the reference file: \texttt{gridInqPosition}}
\index{gridInqPosition}
\label{gridInqPosition}

The function {\texttt{gridInqPosition}} returns the position of grid in the reference file.

\subsubsection*{Usage}

\begin{verbatim}
    INTEGER FUNCTION gridInqPosition(INTEGER gridID)
\end{verbatim}

\hspace*{4mm}\begin{minipage}[]{15cm}
\begin{deflist}{\texttt{gridID}\ }
\item[\texttt{gridID}]
Grid ID, from a previous call to {\htmlref{\texttt{gridCreate}}{gridCreate}} or {\htmlref{\texttt{vlistInqVarGrid}}{vlistInqVarGrid}}.

\end{deflist}
\end{minipage}

\subsubsection*{Result}

{\texttt{gridInqPosition}} returns the position of grid in the reference file.


\subsection{Define the reference URI for an unstructured grid: \texttt{gridDefReference}}
\index{gridDefReference}
\label{gridDefReference}

The function {\texttt{gridDefReference}} defines the reference URI for an unstructured grid.

\subsubsection*{Usage}

\begin{verbatim}
    SUBROUTINE gridDefReference(INTEGER gridID, CHARACTER*(*) reference)
\end{verbatim}

\hspace*{4mm}\begin{minipage}[]{15cm}
\begin{deflist}{\texttt{reference}\ }
\item[\texttt{gridID}]
Grid ID, from a previous call to {\htmlref{\texttt{gridCreate}}{gridCreate}}.
\item[\texttt{reference}]
Reference URI for an unstructured grid.

\end{deflist}
\end{minipage}


\subsection{Get the reference URI to an unstructured grid: \texttt{gridInqReference}}
\index{gridInqReference}
\label{gridInqReference}

The function {\texttt{gridInqReference}} returns the reference URI to an unstructured grid.

\subsubsection*{Usage}

\begin{verbatim}
    char *gridInqReference(INTEGER gridID, CHARACTER*(*) reference)
\end{verbatim}

\hspace*{4mm}\begin{minipage}[]{15cm}
\begin{deflist}{\texttt{gridID}\ }
\item[\texttt{gridID}]
Grid ID, from a previous call to {\htmlref{\texttt{gridCreate}}{gridCreate}} or {\htmlref{\texttt{vlistInqVarGrid}}{vlistInqVarGrid}}.

\end{deflist}
\end{minipage}

\subsubsection*{Result}

{\texttt{gridInqReference}} returns the reference URI to an unstructured grid.


\subsection{Define the UUID for an unstructured grid: \texttt{gridDefUUID}}
\index{gridDefUUID}
\label{gridDefUUID}

The function {\texttt{gridDefUUID}} defines the UUID for an unstructured grid.

\subsubsection*{Usage}

\begin{verbatim}
    SUBROUTINE gridDefUUID(INTEGER gridID, CHARACTER*(*) uuid)
\end{verbatim}

\hspace*{4mm}\begin{minipage}[]{15cm}
\begin{deflist}{\texttt{gridID}\ }
\item[\texttt{gridID}]
Grid ID, from a previous call to {\htmlref{\texttt{gridCreate}}{gridCreate}}.
\item[\texttt{uuid}]
UUID for an unstructured grid.

\end{deflist}
\end{minipage}


\subsection{Get the UUID to an unstructured grid: \texttt{gridInqUUID}}
\index{gridInqUUID}
\label{gridInqUUID}

The function {\texttt{gridInqUUID}} returns the UUID to an unstructured grid.

\subsubsection*{Usage}

\begin{verbatim}
    SUBROUTINE gridInqUUID(INTEGER gridID, CHARACTER*(*) uuid)
\end{verbatim}

\hspace*{4mm}\begin{minipage}[]{15cm}
\begin{deflist}{\texttt{gridID}\ }
\item[\texttt{gridID}]
Grid ID, from a previous call to {\htmlref{\texttt{gridCreate}}{gridCreate}} or {\htmlref{\texttt{vlistInqVarGrid}}{vlistInqVarGrid}}.

\end{deflist}
\end{minipage}

\subsubsection*{Result}

{\texttt{gridInqUUID}} returns the UUID to an unstructured grid to the parameter uuid.
